%% bare_conf.tex
%% V1.4b
%% 2015/08/26
%% by Michael Shell
%% See:
%% http://www.michaelshell.org/
%% for current contact information.
%%
%% This is a skeleton file demonstrating the use of IEEEtran.cls
%% (requires IEEEtran.cls version 1.8b or later) with an IEEE
%% conference paper.
%%
%% Support sites:
%% http://www.michaelshell.org/tex/ieeetran/
%% http://www.ctan.org/pkg/ieeetran
%% and
%% http://www.ieee.org/

%%*************************************************************************
%% Legal Notice:
%% This code is offered as-is without any warranty either expressed or
%% implied; without even the implied warranty of MERCHANTABILITY or
%% FITNESS FOR A PARTICULAR PURPOSE!
%% User assumes all risk.
%% In no event shall the IEEE or any contributor to this code be liable for
%% any damages or losses, including, but not limited to, incidental,
%% consequential, or any other damages, resulting from the use or misuse
%% of any information contained here.
%%
%% All comments are the opinions of their respective authors and are not
%% necessarily endorsed by the IEEE.
%%
%% This work is distributed under the LaTeX Project Public License (LPPL)
%% ( http://www.latex-project.org/ ) version 1.3, and may be freely used,
%% distributed and modified. A copy of the LPPL, version 1.3, is included
%% in the base LaTeX documentation of all distributions of LaTeX released
%% 2003/12/01 or later.
%% Retain all contribution notices and credits.
%% ** Modified files should be clearly indicated as such, including  **
%% ** renaming them and changing author support contact information. **
%%*************************************************************************


% *** Authors should verify (and, if needed, correct) their LaTeX system  ***
% *** with the testflow diagnostic prior to trusting their LaTeX platform ***
% *** with production work. The IEEE's font choices and paper sizes can   ***
% *** trigger bugs that do not appear when using other class files.       ***                          ***
% The testflow support page is at:
% http://www.michaelshell.org/tex/testflow/



\documentclass[conference]{IEEEtran}
% Some Computer Society conferences also require the compsoc mode option,
% but others use the standard conference format.
%
% If IEEEtran.cls has not been installed into the LaTeX system files,
% manually specify the path to it like:
% \documentclass[conference]{../sty/IEEEtran}





% Some very useful LaTeX packages include:
% (uncomment the ones you want to load)
\usepackage{booktabs}

% *** MISC UTILITY PACKAGES ***
%
%\usepackage{ifpdf}
% Heiko Oberdiek's ifpdf.sty is very useful if you need conditional
% compilation based on whether the output is pdf or dvi.
% usage:
% \ifpdf
%   % pdf code
% \else
%   % dvi code
% \fi
% The latest version of ifpdf.sty can be obtained from:
% http://www.ctan.org/pkg/ifpdf
% Also, note that IEEEtran.cls V1.7 and later provides a builtin
% \ifCLASSINFOpdf conditional that works the same way.
% When switching from latex to pdflatex and vice-versa, the compiler may
% have to be run twice to clear warning/error messages.






% *** CITATION PACKAGES ***
%
%\usepackage{cite}
% cite.sty was written by Donald Arseneau
% V1.6 and later of IEEEtran pre-defines the format of the cite.sty package
% \cite{} output to follow that of the IEEE. Loading the cite package will
% result in citation numbers being automatically sorted and properly
% "compressed/ranged". e.g., [1], [9], [2], [7], [5], [6] without using
% cite.sty will become [1], [2], [5]--[7], [9] using cite.sty. cite.sty's
% \cite will automatically add leading space, if needed. Use cite.sty's
% noadjust option (cite.sty V3.8 and later) if you want to turn this off
% such as if a citation ever needs to be enclosed in parenthesis.
% cite.sty is already installed on most LaTeX systems. Be sure and use
% version 5.0 (2009-03-20) and later if using hyperref.sty.
% The latest version can be obtained at:
% http://www.ctan.org/pkg/cite
% The documentation is contained in the cite.sty file itself.






% *** GRAPHICS RELATED PACKAGES ***
%
\ifCLASSINFOpdf
  % \usepackage[pdftex]{graphicx}
  % declare the path(s) where your graphic files are
  % \graphicspath{{../pdf/}{../jpeg/}}
  % and their extensions so you won't have to specify these with
  % every instance of \includegraphics
  % \DeclareGraphicsExtensions{.pdf,.jpeg,.png}
\else
  % or other class option (dvipsone, dvipdf, if not using dvips). graphicx
  % will default to the driver specified in the system graphics.cfg if no
  % driver is specified.
  % \usepackage[dvips]{graphicx}
  % declare the path(s) where your graphic files are
  % \graphicspath{{../eps/}}
  % and their extensions so you won't have to specify these with
  % every instance of \includegraphics
  % \DeclareGraphicsExtensions{.eps}
\fi
% graphicx was written by David Carlisle and Sebastian Rahtz. It is
% required if you want graphics, photos, etc. graphicx.sty is already
% installed on most LaTeX systems. The latest version and documentation
% can be obtained at:
% http://www.ctan.org/pkg/graphicx
% Another good source of documentation is "Using Imported Graphics in
% LaTeX2e" by Keith Reckdahl which can be found at:
% http://www.ctan.org/pkg/epslatex
%
% latex, and pdflatex in dvi mode, support graphics in encapsulated
% postscript (.eps) format. pdflatex in pdf mode supports graphics
% in .pdf, .jpeg, .png and .mps (metapost) formats. Users should ensure
% that all non-photo figures use a vector format (.eps, .pdf, .mps) and
% not a bitmapped formats (.jpeg, .png). The IEEE frowns on bitmapped formats
% which can result in "jaggedy"/blurry rendering of lines and letters as
% well as large increases in file sizes.
%
% You can find documentation about the pdfTeX application at:
% http://www.tug.org/applications/pdftex





% *** MATH PACKAGES ***
%
\usepackage{amsmath}
% A popular package from the American Mathematical Society that provides
% many useful and powerful commands for dealing with mathematics.
%
% Note that the amsmath package sets \interdisplaylinepenalty to 10000
% thus preventing page breaks from occurring within multiline equations. Use:
%\interdisplaylinepenalty=2500
% after loading amsmath to restore such page breaks as IEEEtran.cls normally
% does. amsmath.sty is already installed on most LaTeX systems. The latest
% version and documentation can be obtained at:
% http://www.ctan.org/pkg/amsmath





% *** SPECIALIZED LIST PACKAGES ***
%
%\usepackage{algorithmic}
% algorithmic.sty was written by Peter Williams and Rogerio Brito.
% This package provides an algorithmic environment fo describing algorithms.
% You can use the algorithmic environment in-text or within a figure
% environment to provide for a floating algorithm. Do NOT use the algorithm
% floating environment provided by algorithm.sty (by the same authors) or
% algorithm2e.sty (by Christophe Fiorio) as the IEEE does not use dedicated
% algorithm float types and packages that provide these will not provide
% correct IEEE style captions. The latest version and documentation of
% algorithmic.sty can be obtained at:
% http://www.ctan.org/pkg/algorithms
% Also of interest may be the (relatively newer and more customizable)
% algorithmicx.sty package by Szasz Janos:
% http://www.ctan.org/pkg/algorithmicx




% *** ALIGNMENT PACKAGES ***
%
%\usepackage{array}
% Frank Mittelbach's and David Carlisle's array.sty patches and improves
% the standard LaTeX2e array and tabular environments to provide better
% appearance and additional user controls. As the default LaTeX2e table
% generation code is lacking to the point of almost being broken with
% rspecialtyCodet to the quality of the end results, all users are strongly
% advised to use an enhanced (at the very least that provided by array.sty)
% set of table tools. array.sty is already installed on most systems. The
% latest version and documentation can be obtained at:
% http://www.ctan.org/pkg/array


% IEEEtran contains the IEEEeqnarray family of commands that can be used to
% generate multiline equations as well as matrices, tables, etc., of high
% quality.




% *** SUBFIGURE PACKAGES ***
%\ifCLASSOPTIONcompsoc
%  \usepackage[caption=false,font=normalsize,labelfont=sf,textfont=sf]{subfig}
%\else
%  \usepackage[caption=false,font=footnotesize]{subfig}
%\fi
% subfig.sty, written by Steven Douglas Cochran, is the modern replacement
% for subfigure.sty, the latter of which is no longer maintained and is
% incompatible with some LaTeX packages including fixltx2e. However,
% subfig.sty requires and automatically loads Axel Sommerfeldt's caption.sty
% which will override IEEEtran.cls' handling of captions and this will result
% in non-IEEE style figure/table captions. To prevent this problem, be sure
% and invoke subfig.sty's "caption=false" package option (available since
% subfig.sty version 1.3, 2005/06/28) as this is will preserve IEEEtran.cls
% handling of captions.
% Note that the Computer Society format requires a larger sans serif font
% than the serif footnote size font used in traditional IEEE formatting
% and thus the need to invoke different subfig.sty package options depending
% on whether compsoc mode has been enabled.
%
% The latest version and documentation of subfig.sty can be obtained at:
% http://www.ctan.org/pkg/subfig




% *** FLOAT PACKAGES ***
%
%\usepackage{fixltx2e}
\usepackage{url}
\usepackage{graphicx}

% fixltx2e, the successor to the earlier fix2col.sty, was written by
% Frank Mittelbach and David Carlisle. This package corrects a few problems
% in the LaTeX2e kernel, the most notable of which is that in current
% LaTeX2e releases, the ordering of single and double column floats is not
% guaranteed to be preserved. Thus, an unpatched LaTeX2e can allow a
% single column figure to be placed prior to an earlier double column
% figure.
% Be aware that LaTeX2e kernels dated 2015 and later have fixltx2e.sty's
% corrections already built into the system in which case a warning will
% be issued if an attempt is made to load fixltx2e.sty as it is no longer
% needed.
% The latest version and documentation can be found at:
% http://www.ctan.org/pkg/fixltx2e


%\usepackage{stfloats}
% stfloats.sty was written by Sigitas Tolusis. This package gives LaTeX2e
% the ability to do double column floats at the bottom of the page as well
% as the top. (e.g., "\begin{figure*}[!b]" is not normally possible in
% LaTeX2e). It also provides a command:
%\fnbelowfloat
% to enable the placement of footnotes below bottom floats (the standard
% LaTeX2e kernel puts them above bottom floats). This is an invasive package
% which rewrites many portions of the LaTeX2e float routines. It may not work
% with other packages that modify the LaTeX2e float routines. The latest
% version and documentation can be obtained at:
% http://www.ctan.org/pkg/stfloats
% Do not use the stfloats baselinefloat ability as the IEEE does not allow
% \baselineskip to stretch. Authors submitting work to the IEEE should note
% that the IEEE rarely uses double column equations and that authors should try
% to avoid such use. Do not be tempted to use the cuted.sty or midfloat.sty
% packages (also by Sigitas Tolusis) as the IEEE does not format its papers in
% such ways.
% Do not attempt to use stfloats with fixltx2e as they are incompatible.
% Instead, use Morten Hogholm'a dblfloatfix which combines the features
% of both fixltx2e and stfloats:
%
% \usepackage{dblfloatfix}
% The latest version can be found at:
% http://www.ctan.org/pkg/dblfloatfix




% *** PDF, URL AND HYPERLINK PACKAGES ***
%
%\usepackage{url}
% url.sty was written by Donald Arseneau. It provides better support for
% handling and breaking URLs. url.sty is already installed on most LaTeX
% systems. The latest version and documentation can be obtained at:
% http://www.ctan.org/pkg/url
% Basically, \url{my_url_here}.




% *** Do not adjust lengths that control margins, column widths, etc. ***
% *** Do not use packages that alter fonts (such as pslatex).         ***
% There should be no need to do such things with IEEEtran.cls V1.6 and later.
% (Unless specifically asked to do so by the journal or conference you plan
% to submit to, of course. )


% correct bad hyphenation here
\hyphenation{op-tical net-works semi-conduc-tor}


\begin{document}
%
% paper title
% Titles are generally capitalized except for words such as a, an, and, as,
% at, but, by, for, in, nor, of, on, or, the, to and up, which are usually
% not capitalized unless they are the first or last word of the title.
% Linebreaks \\ can be used within to get better formatting as desired.
% Do not put math or special symbols in the title.
\title{A Conversational Agent for Natural Language Access to Public Health Data}

% Pensei em trocar o título, mas as alternativas abaixo acabaram deixando o título muito longo... Se tiverem alguma outra ideia de título coloquem aqui.
% - Bridging Natural Language and Public Health Databases: A Conversational Agent for DATASUS Analytics
% - Enabling Natural Language Access to DATASUS: A Conversational Agent for Public Health Data Analytics
% Sugestão da Isadora que alterei um pouco:
% - Natural Language Access to Public Health Data: A Conversational Agent
% Como estava antes:
% - A Conversational Agent for Public \\Health Data Analytics


% author names and affiliations
% use a multiple column layout for up to three different
% affiliations
\author{\IEEEauthorblockN{Maicon Moraes}
\IEEEauthorblockA{UFRGS - Federal University of \\Rio Grande do Sul\\
Porto Alegre, Brazil\\
maicon.kevyn@ufrgs.br}
\and
\IEEEauthorblockN{Isadora Figueiredo, Victória Marques, Duncan Ruiz, and \\Isabel H. Manssour}
\IEEEauthorblockA{PUCRS - Pontifical Catholic University of Rio Grande do Sul\\
Porto Alegre, Brazil\\
isabel.manssour@pucrs.br}}

% conference papers do not typically use \thanks and this command
% is locked out in conference mode. If really needed, such as for
% the acknowledgment of grants, issue a \IEEEoverridecommandlockouts
% after \documentclass

% for over three affiliations, or if they all won't fit within the width
% of the page, use this alternative format:
%
%\author{\IEEEauthorblockN{Michael Shell\IEEEauthorrefmark{1},
%Homer Simpson\IEEEauthorrefmark{2},
%James Kirk\IEEEauthorrefmark{3},
%Montgomery Scott\IEEEauthorrefmark{3} and
%Eldon Tyrell\IEEEauthorrefmark{4}}
%\IEEEauthorblockA{\IEEEauthorrefmark{1}School of Electrical and Computer Engineering\\
%Georgia Institute of Technology,
%Atlanta, Georgia 30332--0250\\ Email: see http://www.michaelshell.org/contact.html}
%\IEEEauthorblockA{\IEEEauthorrefmark{2}Twentieth Century Fox, Springfield, USA\\
%Email: homer@thesimpsons.com}
%\IEEEauthorblockA{\IEEEauthorrefmark{3}Starfleet Academy, San Francisco, California 96678-2391\\
%Telephone: (800) 555--1212, Fax: (888) 555--1212}
%\IEEEauthorblockA{\IEEEauthorrefmark{4}Tyrell Inc., 123 Replicant Street, Los Angeles, California 90210--4321}}




% use for special paper notices
%\IEEspecialtyCodeialpapernotice{(Invited Paper)}




% make the title area
\maketitle

% As a general rule, do not put math, special symbols or citations
% in the abstract
\begin{abstract}
%We present the first Text-to-SQL system for DATASUS SIH-RD, Brazil's national hospital information system covering 18.7~million hospitalization records.Despite being one of the world's largest epidemiological repositories, SIH-RD microdata remains analytically inaccessible to non-technical practitioners: opaque clinical encodings, ambiguous join paths, and coded value mappings confound general-purpose language models, and no Portuguese natural-language interface for DATASUS currently exists.
DATASUS, Brazil's national public health data repository, provides access to large volumes of epidemiological data. Among its systems, the Hospital Information System in Reduced Data format (SIH-RD) records millions of hospitalization procedures. Despite being one of the world's largest epidemiological repositories, SIH-RD microdata remains analytically inaccessible to non-technical practitioners: opaque clinical encodings, ambiguous join paths, and coded value mappings confound general-purpose language models, and no Portuguese natural-language interface for DATASUS currently exists. To the best of our knowledge, we present the first Text-to-SQL system for DATASUS SIH-RD, enabling queries over 18.7~million hospitalization records.
To address this, we derive fifteen domain-specific SQL generation rules from systematic SIH-RD schema analysis and embed them in a 9-stage LangGraph pipeline with query routing, automatic table selection, chain-of-thought planning, SQL generation, static validation, and bounded self-repair, requiring no model fine-tuning.
We also construct a benchmark of 120 Portuguese healthcare queries stratified into Easy, Medium, and Hard tiers (40~each) with gold-standard SQL over records from two Brazilian states (2008--2023), comprising the first formal Text-to-SQL evaluation on SIH-RD.
The agent achieves 93.3\% execution accuracy (112/120) with 100\% pipeline completion; a controlled single-shot baseline sharing identical model, domain rules, and prompts achieves 90.0\% (108/120), with the advantage concentrated in Hard queries ($+$10.0~pp), isolating the contribution of graph orchestration for complex multi-table reasoning.
%The domain rules, pipeline, and benchmark are publicly available.
\end{abstract}

% no keywords




% For peer review papers, you can put extra information on the cover
% page as needed:
% \ifCLASSOPTIONpeerreview
% \begin{center} \bfseries EDICS Category: 3-BBND \end{center}
% \fi
%
% For peerreview papers, this IEEEtran command inserts a page break and
% creates the second title. It will be ignored for other modes.
\IEEEpeerreviewmaketitle


%%%%%%%%%%%%%%%%%%%%%%%%%%%%%%%%%%%%%%%%%%%%%%%%%%%%%%%%%%%%%%%%%%%%%%
\section{Introduction}
% no \IEEEPARstart


The evolution of health information technologies has consolidated the availability of large volumes of data, transforming them into strategic resources for health promotion and evidence-based decision-making \cite{Raghupathi2014}. The analysis of clinical data enables the optimization of resource management and the improvement of diagnostic accuracy \cite{Alhejaily2025}. However, knowledge extraction depends on the development of systems, often with complex interfaces, that allow for the dynamic exploration of information \cite{Dash2019}.

The Brazilian Unified Health System (SUS) exemplifies this issue. As a universal, comprehensive, and free system, SUS assists the entire population, including foreigners, and its data constitutes one of the largest epidemiological repositories in the world: DATASUS\footnote{\url{https://datasus.saude.gov.br/}}, with approximately 3 billion procedures recorded in 2023 alone~\cite{Brasil2024}. However, the analytical use of this data is limited by the fragmentation of its multiple information systems, among them the Hospital Information System – Reduced Data (SIH-RD), and by excessive use of coding. Also, sparse documentation and complex schemas, spanning various technical domains such as the International Classification of Diseases (ICD-10) and the National Registry of Health Establishments, hinder interpretation and analysis.
%The Brazilian Unified Health System (SUS) exemplifies this issue. As a universal, comprehensive, and free system, SUS provides assistance to the entire population, including foreigners, and its data constitutes one of the largest epidemiological repositories in the world: DATASUS\footnote{\url{https://datasus.saude.gov.br/}}, with approximately 3 billion procedures recorded in 2023 alone~\cite{Brasil2024}. However, the analytical utilization of this data is limited by the fragmentation of its multiple information systems, as the Hospital Information System – Reduced Data (SIH-RD), and the excessive use of coding. Furthermore, sparse documentation and complex schemas, spanning various technical domains such as the International Classification of Diseases (ICD-10) and the National Registry of Health Establishments, hinder interpretation and analysis.

To mitigate the gap between raw data and end users, Large Language Models (LLMs) are emerging as a promising alternative, enabling the development of natural language interfaces for structured data. Through Text-to-SQL systems, these technologies seek to translate free-form questions into executable queries, reducing the dependence on technical intermediaries~\cite{hong2024survey}. However, although benchmarks such as Spider~\cite{yu2018spider} and BIRD~\cite{li2023bird} demonstrate success in controlled environments, the direct application of LLMs in real-world health ecosystems is still limited. The semantic complexity and fragmentation of databases such as those from DATASUS require more than a direct translation; they demand mechanisms that ensure reasoning integrity and the correct interpretation of specialized domains~\cite{Alhejaily2025}.

In response, this work presents a Portuguese conversational Text-to-SQL agent for DATASUS SIH-RD analytics. The distinguishing contribution is \emph{domain knowledge engineering}: fifteen SQL generation rules and comprehensive SUS-aware value mappings, both derived from systematic SIH-RD schema analysis, that encode the clinical distinctions a general-purpose LLM cannot reliably infer. These are embedded in a 9-stage stateful LangGraph pipeline combining query routing, automatic table selection, chain-of-thought planning for complex queries, domain-aware SQL generation, semantic validation, and bounded self-repair, which ensures 100\% pipeline completion and multi-turn conversational state --- structural properties that follow from graph orchestration but are not assessed by the single-query EX benchmark used here. The pipeline operates over approximately 18 million SIH-RD records structured under a Snowflake schema, without any model fine-tuning. The main contribution of this work are:

\begin{itemize}
\item The presentation of the first NL-to-SQL agent for DATASUS SIH-RD, achieving 93.3\% execution accuracy and 100\% pipeline completion on 120 stratified Portuguese healthcare queries over 18~million hospitalization records.
%\item We present the first NL-to-SQL agent for DATASUS SIH-RD, achieving 93.3\% execution accuracy and 100\% pipeline completion on 120 stratified Portuguese healthcare queries over 18~million hospitalization records.

\item The creation, to our knowledge, of the first formal Text-to-SQL benchmark for DATASUS SIH-RD: 120 queries with gold-standard SQL spanning three difficulty tiers.
%\item We introduce, to our knowledge, the first formal Text-to-SQL benchmark for DATASUS SIH-RD: 120 Portuguese queries with gold-standard SQL spanning three difficulty tiers.

\item A controlled comparison isolating architecture from domain knowledge, showing a directional $+$10.0~pp advantage on Hard queries (McNemar $p{=}0.289$, underpowered at $m{=}8$ discordant pairs), where multi-table joins, and complex aggregations demand the structured reasoning that graph orchestration provides, motivating a larger-scale evaluation.
%\item Through a controlled comparison isolating architecture from domain knowledge, we show a directional $+$10.0~pp advantage on Hard queries (McNemar $p{=}0.289$, underpowered at $m{=}8$ discordant pairs), where multi-table joins and complex aggregations demand the structured reasoning that graph orchestration provides, motivating a larger-scale evaluation.

\end{itemize}

% objetivo

%%%%%%%%%%%%%%%%%%%%%%%%%%%%%%%%%%%%%%%%%%%%%%%%%%%%%%%%%%%%%%%%%%%%%%
\section{Related Work}
\label{sec:related}

An early systematic review found that healthcare chatbots were largely rule-based, with immature clinical evaluations~\cite{laranjo2018}. With LLMs, agentic frameworks now orchestrate external tools and personalize interactions \cite{opencha2025}. In data-intensive contexts, agents that \emph{generate and execute code} demonstrate advantages for complex multi-table reasoning. EHRAgent~\cite{ehragent2024}, e.g., iteratively writes, runs, and debugs code to answer physician queries, outperforming strong baselines on real-world EHR datasets.

These developments align with broader agentic patterns such as planning-and-acting~\cite{react2023}, self-reflection~\cite{reflexion2023}, and multi-agent collaboration~\cite{autogen2023,camel2023}, which are particularly relevant when NL queries must be transformed into faithful, auditable analytics over structured healthcare data. Complementary work explores hybrid access to structured and unstructured sources~\cite{suql2024}. Query decomposition approaches such as DIN-SQL~\cite{pourreza2024din} further improved complex multi-join query handling through sub-question chaining. Together, these advances motivate agentic systems that move beyond informational dialogue to deliver trustworthy, data-driven analytics.

In the Brazilian Unified Healthcare System (SUS), conversational agents have primarily focused on \emph{service guidance and information} rather than analytical querying. Official initiatives include vaccination assistance via WhatsApp~\cite{ms_vax_chatbot_2023} and the Ombudsman service~\cite{ouvsus_chatbot_2024}. Academic research similarly emphasizes informational support and usability~\cite{nascimento2024chatbots,silva2024chatbots}. Prototypes developed for primary care, such as chatbots tested with Family Health Units in Pernambuco~(northeastern state)~\cite{jhi_aps_chatbot_2024}, explore how conversational agents can assist patients and professionals in accessing health information, scheduling appointments, and clarifying doubts. Other initiatives focus on maternal and child health education \cite{xavier2021chatbot}. However, these systems do not incorporate agentic reasoning or access to DATASUS microdata for analytical purposes.

To our knowledge, no conversational agents plan, generate, and validate SQL against DATASUS microdata (e.g., SIH-RD). Brazilian systems focus on informational services~\cite{nascimento2024chatbots,silva2024chatbots}, while international Text-to-SQL work (EHRSQL~\cite{lee2023ehrsql}, EHRAgent~\cite{ehragent2024}) targets well-documented clinical environments. We address this gap by presenting a Portuguese NL$\rightarrow$SQL conversational agent for DATASUS, described next.
%, built over a structured health data warehouse, with schema-aware prompting and a validation-and-repair pipeline, evaluated on a domain-specific benchmark of 120 Portuguese healthcare queries.


%%%%%%%%%%%%%%%%%%%%%%%%%%%%%%%%%%%%%%%%%%%%%%%%%%%%%%%%%%%%%%%%%%%%%%
\section{Proposed Approach}
\label{sec:approach}


\subsection{Data Gathering and Preprocessing}

%Inicialmente, coletamos 18.671.329.658 registros do SIH-RD por meio da biblioteca PySUS\footnote{\url{https://pysus.readthedocs.io/}}, abrangendo dois estados do Brasil (RS e MA) no período de 2008 a 2023. Cada linha vincula um procedimento a uma Autorização de Internação Hospitalar (AIH). Optou-se pelos dados reduzidos, na qual o DATASUS realiza a contração dos registros, para viabilizar o processamento em hardware local convencional, visto que a versão completa exigiria uma grande capacidade computacional.
Initially, 18,671,329 records were collected from the SIH-RD using the PySUS library\footnote{\url{https://pysus.readthedocs.io/}}, covering the states of Rio Grande do Sul and Maranhão from 2008 to 2023, selected as a case study. Each row links a procedure to a Hospital Admission Authorization (AIH). The reduced data format (RD) was chosen because DATASUS performs record contraction in this version, enabling processing on commodity local hardware, whereas the full version would require excessive computational capacity. Furthermore, the dataset comprises public microdata containing no personally identifiable information (PII), ensuring compliance with the Brazilian General Data Protection Law (LGPD) and ethical standards for public-domain research.

%A manipulação foi executada com a biblioteca \textbf{Polars}\footnote{\url{https://pola.rs/}}, utilizando Lazy Evaluation e processamento em blocos para otimizar a memória RAM~\cite{2024polars}. O pré-processamento estruturou-se em três frentes: (1) \textbf{Otimização de Tipagem}, via downcasting sistemático (e.g., $Int16 \rightarrow Int8$) e tratamento de nulos; (2) \textbf{Integridade de Regras de Negócio}, com a derivação da variável de permanência (diferença entre as datas de saída e internação) e cálculo de idade;  (3) \textbf{Saneamento de Strings}, com padronização de campos de diagnóstico (CID-10) e de municípios.  Além de outras métricas operacionais
The preprocessing was aligned with data quality and management principles defined in the DAMA-DMBOK~\cite{dama2017}, and was structured across three fronts: (1) \textbf{Type Optimization}, via systematic downcasting (e.g., $Int16 \rightarrow Int8$) and null handling; (2) \textbf{Business Rule Integrity}, with the derivation of the length of stay variable, calculated as the difference between discharge and admission dates, and age calculation; (3) \textbf{String Sanitization}, including the standardization of diagnostic fields (ICD-10) and municipalities, and other operational metrics.




\subsection{Database Modeling and Loading}

%O banco de dados foi estruturado no modelo Snowflake, otimizado para ambientes OLAP, e conta com duas tabelas de fatos centrais: atendimentos e internações. Para mitigar a complexidade da codificação nativa do DATASUS e aumentar a interpretabilidade da LLM, foram implementadas dimensões normalizadas que traduzem códigos técnicos em descritores semânticos. A carga dos dados foi executada via DuckDB\footnote{\url{https://duckdb.org/}}, aproveitando a integração nativa com a biblioteca Polars para realizar bulk inserts, com um mapeamento rigoroso de tipos primitivos. Essa arquitetura, consolidada por chaves primárias e estrangeiras, assegura a integridade referencial necessária para operações complexas de junção e de agregação analítica.
The database was structured in a \textbf{Snowflake schema}, a normalized multidimensional structure optimized for OLAP  analytical workloads~\cite{raasveldt2020data}, comprising
one central fact table, one bridge table, and 14 dimension tables~\cite{kimball2013data}. To mitigate the complexity of native DATASUS coding and increase LLM interpretability, normalized dimensions were implemented to translate technical codes into semantic descriptors. Data loading was executed via \textbf{DuckDB}\footnote{\url{https://duckdb.org/}}, leveraging its native integration with the Polars library to perform bulk inserts with rigorous primitive type mapping~\cite{raasveldt2020data}. This architecture, consolidated by primary and foreign keys, ensures the referential integrity required for complex join operations and analytical aggregation~\cite{dama2017}.


\subsection{Conversational Agent}

Our conversational agent translates natural-language Portuguese queries into executable SQL queries against the described structured health data warehouse. Building on stateful graph orchestration (LangGraph\footnote{\url{https://www.langchain.com/langgraph}}), the agent decomposes NL$\rightarrow$SQL into modular steps (routing, schema pruning, chain-of-thought planning for complex queries, SQL synthesis, validation, and execution) while maintaining explicit state for error recovery and auditability. This design follows recent evidence that graph-orchestrated, tool-augmented pipelines with self-correction improve robustness and execution accuracy in Text-to-SQL~\cite{react2023,reflexion2023}.

The agent uses GPT-4o-mini~\cite{openai2024gpt4omini} (via OpenAI API) with temperature=0 throughout, reducing sampling variability in SQL generation but not guaranteeing determinism. Each pipeline stage receives domain-specific metadata (table schemas, column types, value ranges, foreign key relationships) together with curated per-table few-shot examples; for instance, the \textit{Encounter} template covers gender filtering (\texttt{adminSex} = 1 for male), temporal aggregations (\texttt{EXTRACT(YEAR FROM admissionDate)}), and JOINs with dimension tables for human-readable descriptions. No model fine-tuning is required.


% The conversational agent employs an agentic workflow with LangGraph providing stateful orchestration to translate Portuguese queries into SQL, execute them against PostgreSQL, and return natural language responses

% (Figure~\ref{images:flux}).

\subsubsection{Schema-Aware Prompting Rules}

The system prompt encodes fifteen domain-specific SQL generation rules (RULES A--O), each derived from SIH-RD schema analysis and DATASUS domain conventions to address a class of semantic errors that arise when general-purpose models encounter domain-coded health databases~\cite{Alhejaily2025}. Their absolute contribution relative to zero-shot prompting is left as future work; both the agent and the baseline use identical rules, so current evaluation isolates architectural contribution given fixed domain knowledge.


\begin{table}[!t]
\centering
\caption{Domain-specific SQL generation rules (RULES A--O).}
\label{tab:rules}
\renewcommand{\arraystretch}{1.1}
\begin{tabular}{@{}lp{6.2cm}@{}}
\hline
\textbf{Rule} & \textbf{Constraint} \\
\hline
A & UTI queries: \texttt{WHERE "ICUCost" > 0}; obstetric: \texttt{specialtyCode = 2} \\
B & Death-cause queries join \texttt{deathICD}; diagnosis queries join \texttt{primaryICD} \\
C & \texttt{LIMIT} only when question explicitly requests top-\textit{N}; never a default limit \\
D & Apply only filters explicitly mentioned in the question; no implicit \texttt{NULL} or mortality guards \\
E & Include \texttt{ICD} code column only when question explicitly says ``with code'' \\
F & Singular ``which X plus Y'' $\to$ \texttt{LIMIT 1}; plural ``what are the \textit{N} X'' $\to$ \texttt{LIMIT N} \\
G & Date filters use \texttt{EXTRACT} on \texttt{admissionDate}; never a non-equijoin with \texttt{tempo} \\
H & \texttt{birthDate} (int, 0--130) for all age comparisons; \texttt{admissionAge} (date) only for birth-year queries \\
I & \texttt{COUNT(*)} for ``how many records''; \texttt{COUNT(DISTINCT col)} only for unique-value counts \\
J & Per-group top-\textit{N} requires \texttt{ROW\_NUMBER() OVER (PARTITION BY ...)}; plain \texttt{LIMIT} is forbidden \\
K & Absence patterns (``never'', ``without'') use \texttt{NOT EXISTS}; \texttt{NOT IN} is forbidden due to NULL semantics \\
L & Side-by-side comparisons use wide-format \texttt{CASE WHEN} pivot; long format is forbidden \\
M & Aggregate on the fact table first, then \texttt{JOIN} for labels; never \texttt{JOIN}-then-group \\
N & Global/local average: compute reference aggregate over the full dataset scope, never avg-of-avgs \\
O & Human-readable output: always \texttt{JOIN} lookup table for descriptions; never return raw codes \\
\hline
\end{tabular}
\end{table}

Table~\ref{tab:rules} presents all rules; RULES B and H are highlighted below to illustrate the depth of required domain encoding.
For instance, RULE B mandates distinct join columns for death-cause queries (\textit{``causa da morte''}) using \texttt{deathICD} versus admission-diagnosis queries using \texttt{primaryICD}, a distinction GPT-4o-mini systematically confuses without explicit instruction. RULE H enforces the pre-calculated integer \texttt{AdmissionAge} (0--130) for all age comparisons, with \texttt{birthDate} reserved for birth-year queries only; violating this rule produces syntactically valid but semantically incorrect results.

Complementing RULES A--O, the prompt includes SUS-specific value mappings
(\texttt{adminSex}: 1=Male, 3=Female; \texttt{deathIndicator}: boolean; \texttt{specialtyCode} codes) and a long-format guard for the \texttt{socioeconomicIndicator} table (mandatory \texttt{WHERE metrica=?} filter to prevent full-table scans). All mappings are centralized in a dedicated schema-enrichment component and injected at SQL generation time, ensuring consistency across the pipeline without duplicating domain knowledge in each prompt template.

As illustrated in Figure~\ref{images:flux5.png}, the pipeline comprises nine stages for DATABASE queries:
\\(1)~\textbf{Query classification} routes the input; DATABASE queries enter the SQL pipeline; conversational, schema, and ambiguous queries are handled directly without SQL generation.
\\(2)~\textbf{Table selection} identifies the relevant table subset using a keyword-regex fast-path (high-confidence matches skip the LLM call) or an LLM-based selector for ambiguous queries, narrowing the schema context passed to subsequent stages and reducing LLM exposure to irrelevant table structures.
\\(3)~\textbf{Schema retrieval} fetches column names, types, foreign keys, and value ranges exclusively for the selected tables, providing a focused context that avoids prompt pollution from the full 13-table schema.
\\(4)~\textbf{Chain-of-thought planning}, applied to structurally complex queries (e.g., per-group top-\textit{N}, global-vs-local averages, temporal comparisons), generates a structured SQL sketch identifying required tables, join paths, aggregation patterns, and domain-specific pitfalls before full SQL synthesis. Simpler queries bypass this stage.
\\(5)~\textbf{SQL generation} synthesizes a candidate query using RULES A--O, SUS-specific value mappings, the CoT plan when available, and per-table few-shot examples.
\\(6)~\textbf{SQL validation} performs a static syntax and schema alignment check via LangChain's SQLDatabaseToolkit; invalid SQL is routed to regeneration, avoiding a database round-trip for syntactically invalid candidates and accelerating the repair cycle.
\\(7)~\textbf{SQL execution} runs the validated query against DuckDB~\cite{raasveldt2019duckdb}.
\\(8)~\textbf{Self-repair}, triggered on validation or execution failure, classifies the error as schema-level (rerouting to table selection for revision) or syntax-level (targeted SQL regeneration with verbatim error context), subject to two generation retries, three validation retries, a global 15-step cap, and a loop-detection guard that exits if two attempts produce the same SQL. %exits if two consecutive attempts produce identical SQL.
\\(9)~\textbf{Response generation} formats the query result as a natural-language response and serves as the guaranteed exit point for all error paths, ensuring 100\% pipeline completion regardless of query complexity; the marginal EX contribution of self-repair remains to be quantified through stage-level ablation.





\begin{figure}[!t]
    \centering
    \includegraphics[width=0.70\linewidth]{images/flux5.png} % coloque o nome exato do arquivo
    \caption{LangGraph sequential workflow, showing the interaction between LLM reasoning, schema-aware SQL generation, and DuckDB execution. Gray boxes denote LLM-driven tasks (reasoning and generation), blue boxes represent database operations, and green boxes represent validation steps.}
    \label{images:flux5.png}
\end{figure}


% \subsubsection{Validation and Self-Repair}

% SQL generation does not always produce executable queries on the first attempt under schema
% heterogeneity and domain-coded values~\cite{jiang2023repairing}. Before reaching DuckDB,
% each candidate query passes a pre-execution static check via LangChain's \texttt{sql\_db\_query\_checker}
% tool, which verifies column references against the retrieved schema and routes invalid SQL
% directly to repair without incurring a database round-trip.

% When validation or execution fails, the agent classifies the error and selects one of two
% repair paths. Schema-level errors (e.g., non-existent tables or columns) feed the database error
% message back into the table-selection step so the LLM can revise the table set before
% regenerating SQL. Syntax-level errors (e.g., type mismatches, invalid operators) trigger targeted
% SQL regeneration without re-running table selection, providing the original question, the
% failing query, and the verbatim error as repair context.

% Both paths are subject to bounded retry limits: at most two generation retries, three
% validation retries, and a global 15-step graph-cycle cap, after which the agent exits
% gracefully with an informative response. A loop-detection guard terminates repair early
% if two consecutive attempts produce identical SQL, preventing degenerate regeneration cycles.


%%%%%%%%%%%%%%%%%%%%%%%%%%%%%%%%%%%%%%%%%%%%%%%%%%%%%%%%%%%%%%%%%%%%%%
\section{Evaluation Methodology}
\label{sec:methodology}


We adopt \textbf{Execution Accuracy (EX)} as our primary metric~\cite{yu2018spider,hong2024survey,zhang2023actsql,mohammadjafari2024review}: a predicted query is correct if its execution returns the same result set as the gold-standard SQL, directly measuring user-facing correctness regardless of syntactic form.
%\paragraph{Benchmark.}
We developed a benchmark of 120 NL queries mapped to gold-standard SQL over the SIH-RD/SUS database (2008--2023, Rio Grande do Sul and Maranh\~ao states), manually constructed by the research team. Queries span four analytical dimensions: aggregation, temporal filtering, geographic analysis, and multi-table joins. These dimensions are intentionally overlapping rather than mutually exclusive: 118/120 queries contain explicit aggregation, 25/120 include temporal constraints or calendar extraction, 31/120 involve geographic filtering or grouping, and 55/120 require at least one JOIN. They are stratified into \textbf{40 Easy}, \textbf{40 Medium}, and \textbf{40 Hard} queries based on structural and semantic complexity; the concentration of JOIN-heavy queries rises from 0/40 in Easy to 20/40 in Medium and 35/40 in Hard, while temporal and geographic reasoning also become more frequent in the harder tiers. Importantly, the analytical dimensions do not determine difficulty by themselves: an aggregation-only question can remain Easy, whereas the combination of aggregation with JOINs, temporal comparison, geographic grouping, or multi-step calculations is what typically pushes a query into Medium or Hard. Accordingly, the benchmark should be interpreted as an initial feasibility instrument for this DATASUS setting, not as a definitive validation of real-world coverage or deployment readiness. %Participation of clinical domain specialists in benchmark construction is planned as future work.

%\paragraph{Baseline construction.}
To isolate the net architectural contribution of the LangGraph pipeline, we implemented a single-shot prompt baseline using the same GPT-4o-mini model~\cite{openai2024gpt4omini} (temperature=0) and \emph{identical} domain prompts: the same RULES A--O, SUS-aware value mappings, per-table few-shot examples, and full schema context. Unlike the LangGraph agent, the baseline issues a single LLM call per query with no table selection, validation, or repair stages. Any difference in EX between the two systems, therefore, reflects the net effect of multi-step graph orchestration at the whole-pipeline level, holding all other factors constant. However, this comparison does not isolate the marginal contribution of internal stages such as table selection, chain-of-thought planning, validation, or self-repair individually. This design avoids the confound present in many agent evaluations, where the baseline uses weaker prompts than the agent.

%\paragraph{Statistical comparison.}
Because both systems are evaluated on the same 120 queries, their per-query outcomes are paired and dependent; standard two-proportion z-tests are therefore inappropriate.
We use McNemar's exact test~\cite{dietterich1998} on the $2{\times}2$ contingency table of discordant pairs ($b$ = queries only the agent answers correctly;
$c$ = queries only the baseline answers correctly), computing the two-sided binomial probability  $p = 2\sum_{k=b}^{m}{\binom{m}{k}(0.5)^m}$, where $m{=}b{+}c$ is the total number of discordant pairs. The exact variant is used because high-accuracy systems on small benchmarks yield few discordant pairs, violating the $b{+}c \geq 25$ requirement of the chi-square approximation. Additionally, we report Wilson 95\% {confidence intervals (CIs)} for each system's overall EX proportion independently, computed as $\hat{p} = k/n$ where $k$ is the number of correct queries and $n{=}120$; Wilson CIs are preferred over the normal approximation for proportions near 1.
%\vspace{0.3cm}
Table~\ref{tab:benchmark_tiers} summarizes the stratification criteria and a representative gold-standard SQL query for each tier (queries translated to English).

\begin{table}[!t]
\centering
\caption{Benchmark stratification: structural criteria and representative gold-standard SQL per tier.}
\label{tab:benchmark_tiers}
\renewcommand{\arraystretch}{1.3}
\begin{tabular}{@{}p{0.7cm}cp{2.5cm}p{4.3cm}@{}}
\hline
\textbf{Tier} & \textbf{n} & \textbf{Structural Criteria} & \textbf{Representative SQL} \\
\hline\hline
\rule{0pt}{8pt}%
Easy   & 40 & Single-table; basic filter or aggregation; no joins &
\parbox[t]{4.2cm}{\ttfamily\small%\vspace{2pt}
SELECT COUNT(*)\\
FROM Encounter\vspace{4pt}} \\
\hline
\rule{0pt}{8pt}%
Medium & 40 & Two-table FK join, or single-table with temporal / grouped aggregation &
\parbox[t]{4.2cm}{\ttfamily\small%\vspace{2pt}
SELECT m.name, COUNT(*) AS n\\
FROM Encounter e\\
JOIN Location m\\
\hspace{3mm}ON e.locId = m.code\\
GROUP BY m.name\\
ORDER BY n DESC~LIMIT~5\vspace{4pt}} \\
\hline
\rule{0pt}{8pt}%
Hard   & 40 & $\geq$2-table joins, multi-step reasoning, or computed metrics across dimensional hierarchies &
\parbox[t]{4.2cm}{\ttfamily\small%\vspace{2pt}
SELECT edu.descr,\\
\hspace{3mm}AVG(e.totalCost)~AS~avgCost\\
FROM Encounter e\\
JOIN EducationCode edu\\
\hspace{3mm}ON e.eduCode~=~edu.code\\
WHERE~edu.descr~IS~NOT~NULL\\
GROUP BY edu.descr\\
ORDER BY avgCost DESC\vspace{4pt}} \\
\hline
\end{tabular}
\end{table}



%%%%%%%%%%%%%%%%%%%%%%%%%%%%%%%%%%%%%%%%%%%%%%%%%%%%%%%%%%%%%%%%%%%%%%
\section{Results and Discussion}
\label{sec:results}


%\paragraph{Quantitative results.}
Both systems returned a response for all 120 queries without pipeline crashes or timeouts (completion rate: 100\%). The LangGraph agent achieved 93.3\% EX (112/120), with strong performance across all tiers: Easy 100.0\% (40/40), Medium 97.5\% (39/40), and Hard 82.5\% (33/40). The prompt baseline achieved 90.0\% EX (108/120): Easy 97.5\% (39/40), Medium 100.0\% (40/40), and Hard 72.5\% (29/40), yielding a 3.3~pp overall advantage for the LangGraph agent. However, the paired comparison remains statistically inconclusive because it is driven by only eight discordant queries overall ($b{=}6$, $c{=}2$, $m{=}8$; McNemar's exact test $p{=}0.289$), which leaves limited power to detect moderate effects even with $n{=}120$ total queries; 95\% Wilson CI: agent [87.4\%, 96.6\%], baseline [83.3\%, 94.2\%]. The advantage is concentrated in the Hard tier ($+$10.0~pp), where multi-table joins, per-group aggregations, and global-vs-local comparisons demand the structured reasoning provided by table selection and chain-of-thought planning. Both discordant cases where the baseline succeeds and the agent fails ($c{=}2$, one Medium, one Hard) reflect table-selection errors that propagate semantically plausible but incorrect queries to execution a failure mode that bypasses syntactic repair and represents the pipeline's dominant observed vulnerability.

Two failure patterns illustrate the performance gap. First, the pipeline's six agent-only successes ($b{=}6$, five in Hard) involve queries requiring per-group top-$N$ rankings or global-versus-local average comparisons: the baseline generates a syntactically valid but structurally incorrect query (e.g., a plain \texttt{LIMIT} instead of \texttt{ROW\_NUMBER() OVER (PARTITION BY \ldots)} violating RULE~J), while the agent's chain-of-thought planning stage produces the correct window function. Second, both agent-only failures ($c{=}2$) share a common mode: the table-selection step maps an ambiguous query to a plausible but incorrect table set (e.g., selecting \texttt{primaryICD} for a death-cause query that requires \texttt{deathICD} per RULE~B), generating a semantically wrong yet syntactically valid query that executes without triggering repair. This asymmetry---CoT planning reducing Hard-tier failures while table-selection errors introduce Medium-tier regressions---motivates targeted ablation of each stage.


\begin{table}[t]
\centering
\caption{Execution Accuracy (EX) by system and difficulty tier.
Exact McNemar's test on overall EX: $b{=}6$, $c{=}2$, $p{=}0.289$.}
\label{tab:agent_results}
\begin{tabular}{lcccc}
\toprule
\textbf{System} & \textbf{Easy} & \textbf{Medium} & \textbf{Hard} & \textbf{Overall} \\
                & \textit{(n=40)} & \textit{(n=40)} & \textit{(n=40)} & \textit{(n=120)} \\
\midrule
Prompt Baseline & 97.5\% & 100.0\% & 72.5\% & 90.0\% \\
LangGraph Agent      & 100.00\% & 97.5\% & 82.5\% & \textbf{93.3\%} \\
\midrule
$\Delta$ (Agent $-$ Baseline) & $+$2.5 pp & $-$2.5 pp & $+$10.0 pp & $+$3.3 pp \\
\bottomrule
\end{tabular}
\end{table}


%\paragraph{Interpreting the accuracy gap.}
The 3.3 pp overall advantage of the LangGraph agent is informative in its structure. Both systems share \emph{identical} domain knowledge (RULES A--O, SUS value mappings, per-table few-shot examples), so any difference in EX reflects the net effect of graph orchestration: table selection, chain-of-thought planning, static validation, and bounded self-repair. The six discordant queries where the agent succeeds and the baseline fails ($b{=}6$) concentrate in the Hard tier (5 out of 6), confirming that the pipeline's structured reasoning stages add value precisely where query complexity is highest. Among the observed cases where the baseline succeeds and the agent fails ($c{=}2$, one Medium, one Hard), both involve table-selection errors that propagate semantically plausible but incorrect queries to execution without triggering syntactic repair; a systematic failure analysis remains as future work; by tier, the agent's 8 errors concentrate in Hard (7) and Medium (1), and the baseline's 12 in Hard (11) and Easy (1), with 6 Hard queries that neither system answers correctly representing the benchmark's current difficulty ceiling. The key inferential limitation is not only the benchmark's current size, but the fact that the paired test is driven by just $m{=}b{+}c{=}8$ discordant pairs. Under those conditions, $p{=}0.289$ does not justify a superiority claim, but neither does it provide strong evidence against a moderate orchestration effect concentrated in Hard queries. The Hard-tier gap of $+$10.0~pp is the strongest directional signal and aligns with the design intent of the pipeline: table selection narrows schema context and chain-of-thought planning scaffolds complex multi-step reasoning, both of which matter most for Hard queries. However, the Medium-tier regression ($-$2.5~pp) is also practically important: because the overall gain is small, the current results do not support a claim of uniform superiority over a simpler single-shot approach. If real-world usage is dominated by medium-complexity requests, the added orchestration overhead may not be justified until the table-selection failure mode is reduced.

%\paragraph{Why LangGraph rather than a bare LLM call.}
The stateful LangGraph design also embeds structural characteristics absent from single-shot systems by construction: typed error recovery distinguishing schema-level from syntax-level failures, multi-turn conversational state for iterative query refinement, bounded execution with loop-detection guards, and a per-decision inspection trail supporting auditability in healthcare information systems. These properties were not evaluated by the current benchmark, which measures only single-query EX; their practical value in deployed clinical analytics settings remains to be demonstrated in future work.

%\subsection*{Limitations}

This evaluation has four principal limitations. First, component-level ablation is absent: the 3.3~pp overall advantage cannot be attributed to individual stages, and the table-selection step is simultaneously a strength (schema disambiguation) and a liability (upstream errors that bypass repair). Second, the main inferential bottleneck is the low number of discordant pairs in the paired comparison: $p{=}0.289$ arises from only $m{=}8$ discordant queries, leaving McNemar's test underpowered for moderate effects even though the benchmark contains 120 total queries. Third, the benchmark was constructed entirely by the research team without clinical specialist participation, which may limit query diversity and ecological validity. Fourth, the study does not quantify the operational overhead of graph orchestration relative to the single-shot baseline: additional LLM/tool stages may increase latency, token usage, and deployment complexity, so the practical cost-benefit of the observed 3.3~pp gain remains to be established.


%\subsection*{Future Work}

%\textbf{Benchmark validation by clinical specialists} is the most pressing gap. The current 81-query benchmark was constructed entirely by the research team. Involving epidemiologists and hospital administrators in query design and gold-standard SQL review would improve ecological validity, coverage of clinically relevant scenarios, and the benchmark's utility as a shared evaluation resource.

%\textbf{Systematic error analysis} is the most direct extension of the current results. The agent produced 3 incorrect queries and the baseline 9, yet this paper does not characterize failure modes by tier, query type, or pipeline stage. Identifying whether errors concentrate in table selection, value grounding, or multi-join reasoning would directly inform targeted improvements.

%\textbf{Model comparison}, including open-source alternatives such as LLaMA 3 and Mistral, would assess whether the domain-specific prompting framework transfers across model families and reduce dependence on a proprietary API, an important consideration for privacy-sensitive healthcare deployments.

%\textbf{Fine-tuning versus prompting} is the natural scientific question raised by the main finding: since prompt engineering drives most of the performance, it is worth investigating whether a smaller model fine-tuned on DATASUS-specific SQL examples could match or exceed the current results at lower cost and latency.

%\textbf{Component-level ablation} would quantify each pipeline stage's marginal contribution to EX by disabling table selection, validation, and self-repair individually, complementing the architectural comparison already established between the agent and the single-shot baseline.


%%%%%%%%%%%%%%%%%%%%%%%%%%%%%%%%%%%%%%%%%%%%%%%%%%%%%%%%%%%%%%%%%%%%%%
\section{Conclusion}
\label{sec:conclusions}
We presented a conversational agent that translates Portuguese natural-language queries into executable SQL over SIH-RD/SUS hospital microdata (18 million records), addressing the lack of accessible analytical tools for non-technical healthcare practitioners. The agent is built on a stateful LangGraph pipeline with 15 domain-specific SQL generation rules (RULES A--O), SUS-aware value mappings, iterative validation, and bounded self-repair. On a benchmark of 120 stratified Portuguese healthcare queries, it achieves 93.3\% EX (112/120) with 100\% pipeline completion. A controlled comparison against a single-shot prompt baseline using the same model and prompts (90.0\% EX, 108/120) shows a 3.3~pp overall difference concentrated in the Hard tier ($+$10.0~pp; McNemar $p{=}0.289$ from $m{=}8$ discordant pairs; 95\% Wilson CI agent: [87.4\%, 96.6\%], baseline: [83.3\%, 94.2\%]), directionally consistent with the pipeline's design intent but statistically underpowered for definitive paired inference at the current benchmark scale. These results therefore support initial feasibility rather than definitive validation of the agent for routine real-world healthcare analytics. Beyond execution accuracy, the LangGraph architecture provides multi-turn conversational state, loop-detection guards, and guaranteed pipeline completion structural properties that follow from the design but whose practical value in deployed clinical analytics settings remains to be demonstrated in future work.


The key contributions of this work are: (1)~\textbf{Domain knowledge engineering for DATASUS SIH-RD}: fifteen SQL generation rules and SUS-aware value mappings derived from schema analysis, embedded in a 9-stage LangGraph pipeline that enables natural-language analytical querying over 18 million records without SQL expertise or model fine-tuning; (2)~\textbf{A domain-specific benchmark}: 120 manually crafted Portuguese queries stratified by complexity (Easy/Medium/Hard) with gold-standard SQL, constituting, to our knowledge, the first formal Text-to-SQL evaluation on SIH-RD/SUS microdata; and (3)~\textbf{A controlled evaluation design}: a single-shot baseline sharing identical model, rules, and prompts with the LangGraph agent, partially isolating the net effect of graph orchestration at the system level and yielding a directional 3.3~pp overall advantage concentrated in the Hard tier ($+$10.0~pp), while leaving component-level contributions for future ablation. A companion repository containing the current implementation and benchmark materials is publicly available\footnote{\url{https://github.com/DAVINTLAB/datasus-conversational-agent/}.}, enabling inspection of the artifact and extension to other public health information systems.

Future work should prioritize two experiments first: stage-level ablation to isolate the individual contributions of table selection, chain-of-thought planning, and self-repair (the current design attributes the Hard-tier gap to their combined effect, but cannot disaggregate them); and a simpler internal baseline that removes the current domain-rule scaffolding, to measure how much of the observed performance comes from prompt-based domain engineering versus orchestration. After these two priorities, the next steps are systematic error characterization of the 8 agent errors and 12 baseline errors by stage and query type to directly inform targeted improvements; benchmark validation with clinical specialists to improve validity; and broader Text-to-SQL comparisons, including open-source models (e.g., LLaMA) and alternative system designs, to position the DATASUS results relative to the wider literature and reduce dependence on proprietary APIs in privacy-sensitive healthcare deployments.





% conference papers do not normally have an appendix


% use section* for acknowledgment
\section*{Acknowledgment}
%The authors gratefully acknowledge the financial support of the CNPq Scholarship PIBIC and Technological Development Productivity (303208/2023-6), and the FAPERGS 22/2551-0000390-7 (RITE CIARS Inteligência Artificial Aplicada à Saúde).
The authors acknowledge the financial support of CNPq Scholarship PIBIC and Technological Development Productivity (303208/2023-6), and FAPERGS 22/2551-0000390-7 (RITE CIARS Inteligência Artificial Aplicada à Saúde).
While preparing and revising this manuscript, we used ChatGPT and Grammarly to ensure clarity and grammatical precision. The authors are responsible for the content and technical accuracy.
%The authors are responsible for creating the entire content and ensuring technical accuracy.





% An example of a floating figure using the graphicx package.
% Note that \label must occur AFTER (or within) \caption.
% For figures, \caption should occur after the \includegraphics.
% Note that IEEEtran v1.7 and later has special internal code that
% is designed to preserve the operation of \label within \caption
% even when the captionsoff option is in effect. However, because
% of issues like this, it may be the safest practice to put all your
% \label just after \caption rather than within \caption{}.
%
% Reminder: the "draftcls" or "draftclsnofoot", not "draft", class
% option should be used if it is desired that the figures are to be
% displayed while in draft mode.
%
%\begin{figure}[!t]
%\centering
%\includegraphics[width=2.5in]{myfigure}
% where an .eps filename suffix will be assumed under latex,
% and a .pdf suffix will be assumed for pdflatex; or what has been declared
% via \DeclareGraphicsExtensions.
%\caption{Simulation results for the network.}
%\label{fig_sim}
%\end{figure}

% Note that the IEEE typically puts floats only at the top, even when this
% results in a large percentage of a column being occupied by floats.


% An example of a double column floating figure using two subfigures.
% (The subfig.sty package must be loaded for this to work.)
% The subfigure \label commands are set within each subfloat command,
% and the \label for the overall figure must come after \caption.
% \hfil is used as a separator to get equal spacing.
% Watch out that the combined width of all the subfigures on a
% line do not exceed the text width or a line break will occur.
%
%\begin{figure*}[!t]
%\centering
%\subfloat[Case I]{\includegraphics[width=2.5in]{box}%
%\label{fig_first_case}}
%\hfil
%\subfloat[Case II]{\includegraphics[width=2.5in]{box}%
%\label{fig_second_case}}
%\caption{Simulation results for the network.}
%\label{fig_sim}
%\end{figure*}
%
% Note that often IEEE papers with subfigures do not employ subfigure
% captions (using the optional argument to \subfloat[]), but instead will
% reference/describe all of them (a), (b), etc., within the main caption.
% Be aware that for subfig.sty to generate the (a), (b), etc., subfigure
% labels, the optional argument to \subfloat must be present. If a
% subcaption is not desired, just leave its contents blank,
% e.g., \subfloat[].


% An example of a floating table. Note that, for IEEE style tables, the
% \caption command should come BEFORE the table and, given that table
% captions serve much like titles, are usually capitalized except for words
% such as a, an, and, as, at, but, by, for, in, nor, of, on, or, the, to
% and up, which are usually not capitalized unless they are the first or
% last word of the caption. Table text will default to \footnotesize as
% the IEEE normally uses this smaller font for tables.
% The \label must come after \caption as always.
%
%\begin{table}[!t]
%% increase table row spacing, adjust to taste
%\renewcommand{\arraystretch}{1.3}
% if using array.sty, it might be a good idea to tweak the value of
% \extrarowheight as needed to properly center the text within the cells
%\caption{An Example of a Table}
%\label{table_example}
%\centering
%% Some packages, such as MDW tools, offer better commands for making tables
%% than the plain LaTeX2e tabular which is used here.
%\begin{tabular}{|c||c|}
%\hline
%One & Two\\
%\hline
%Three & Four\\
%\hline
%\end{tabular}
%\end{table}


% Note that the IEEE does not put floats in the very first column
% - or typically anywhere on the first page for that matter. Also,
% in-text middle ("here") positioning is typically not used, but it
% is allowed and encouraged for Computer Society conferences (but
% not Computer Society journals). Most IEEE journals/conferences use
% top floats exclusively.
% Note that, LaTeX2e, unlike IEEE journals/conferences, places
% footnotes above bottom floats. This can be corrected via the
% \fnbelowfloat command of the stfloats package.




% trigger a \newpage just before the given reference
% number - used to balance the columns on the last page
% adjust value as needed - may need to be readjusted if
% the document is modified later
%\IEEEtriggeratref{8}
% The "triggered" command can be changed if desired:
%\IEEEtriggercmd{\enlargethispage{-5in}}

% references section

% can use a bibliography generated by BibTeX as a .bbl file
% BibTeX documentation can be easily obtained at:
% http://mirror.ctan.org/biblio/bibtex/contrib/doc/
% The IEEEtran BibTeX style support page is at:
% http://www.michaelshell.org/tex/ieeetran/bibtex/
%\bibliographystyle{IEEEtran}
% argument is your BibTeX string definitions and bibliography database(s)
%\bibliography{IEEEabrv,../bib/paper}
%
% <OR> manually copy in the resultant .bbl file
% set second argument of \begin to the number of references
% (used to reserve space for the reference number labels box)


\bibliographystyle{IEEEtran}
\bibliography{sample-base}


% that's all folks
\end{document}
